\documentclass[11pt,a4paper]{article}
\usepackage[utf8]{inputenc}
\usepackage[T1]{fontenc}
\usepackage{amsmath,amssymb,amsthm}
\usepackage[margin=2.5cm]{geometry}
\usepackage{hyperref}

\newtheorem{theorem}{Theorem}
\newtheorem{proposition}[theorem]{Proposition}
\newtheorem{corollary}[theorem]{Corollary}
\newtheorem{lemma}[theorem]{Lemma}
\theoremstyle{definition}
\newtheorem{definition}[theorem]{Definition}
\newtheorem{remark}[theorem]{Remark}

\newcommand{\R}{\mathbb{R}}
\newcommand{\W}{\mathcal{W}}
\newcommand{\PS}{\mathcal{P}(S)}
\DeclareMathOperator*{\argmin}{arg\,min}

\title{Projection of a Wasserstein barycenter on a fixed discrete support}
\author{}
\date{}

\begin{document}
\maketitle

\section{Setting}

Let $\mathcal{P}_2(\R)$ be the set of probability measures on $\R$ with finite
second moment. For $\mu\in\mathcal{P}_2(\R)$ let $F_\mu$ be its distribution
function and
\[
  Q_\mu(t)=\inf\{x\in\R : F_\mu(x)\ge t\},\qquad t\in(0,1),
\]
its (left-continuous) quantile function. It is well known that, on the real
line, the $L_2$ Wasserstein distance is the $L_2$ distance between quantile
functions:
\begin{equation}\label{eq:W2}
  W_2^2(\mu,\nu)=\int_0^1\big(Q_\mu(t)-Q_\nu(t)\big)^2\,dt .
\end{equation}
Moreover, a function $Q:(0,1)\to\R$ is the quantile function of some
$\nu\in\mathcal{P}_2(\R)$ if and only if it is non-decreasing,
left-continuous and square integrable; in that case $\nu$ is the image of the
Lebesgue measure $\lambda$ on $(0,1)$ through $Q$, $\nu=Q_{\#}\lambda$.

Let $\mu_1,\dots,\mu_n\in\mathcal{P}_2(\R)$ and $w_1,\dots,w_n\ge 0$ with
$\sum_i w_i=1$. The (Wasserstein) barycenter of the $\mu_i$ is
\[
  \bar\mu=\argmin_{\nu\in\mathcal{P}_2(\R)}\ \sum_{i=1}^n w_i\,W_2^2(\mu_i,\nu),
  \qquad Q_{\bar\mu}=\sum_{i=1}^n w_i\,Q_{\mu_i}.
\]
Even when all the $\mu_i$ are supported on a common finite set (for instance
the categories $\{1,\dots,7\}$ of a Likert scale), $\bar\mu$ is in general
supported on a larger set of non-integer values.

Let $S=\{s_1<s_2<\dots<s_K\}\subset\R$ be a finite set (the \emph{domain} of
the variable) and
\[
  \PS=\{\nu\in\mathcal{P}_2(\R):\ \nu(S)=1\}
\]
the set of distributions supported on $S$. We look for the distribution of
$\PS$ that best represents the barycenter.

\begin{definition}[Nearest-point map]
Let $m_k=(s_k+s_{k+1})/2$, $k=1,\dots,K-1$, be the midpoints of $S$, and set
$m_0=-\infty$, $m_K=+\infty$. The map $T_S:\R\to S$ is
\[
  T_S(x)=s_k \quad\text{if } x\in(m_{k-1},m_k],\qquad k=1,\dots,K .
\]
\end{definition}

$T_S(x)$ is a point of $S$ at minimum distance from $x$: $|x-T_S(x)|=d(x,S)=\min_{s\in S}|x-s|$.
The minimizer is unique except at the midpoints $x=m_k$, which are equidistant
from $s_k$ and $s_{k+1}$ (the definition assigns them to $s_k$). Moreover $T_S$
is non-decreasing and left-continuous.

The \emph{simplified barycenter} is $\tilde\mu=(T_S)_{\#}\bar\mu$, i.e.
\[
  \tilde\mu(\{s_k\})=\bar\mu\big((m_{k-1},m_k]\big),\qquad k=1,\dots,K:
\]
each value $s_k$ receives the mass of the barycenter lying between the two
adjacent midpoints. This is the rule implemented in
\texttt{DDAWASS::simplifyD()} and \texttt{DDAWASS::WD.vec.mean.simplified()}.

\section{Results}

\begin{lemma}[Quantile function of the projection]\label{lem:quantile}
For every $\mu\in\mathcal{P}_2(\R)$,
$Q_{(T_S)_{\#}\mu}=T_S\circ Q_\mu$.
\end{lemma}

\begin{proof}
$T_S\circ Q_\mu$ is non-decreasing and left-continuous, being the composition
of two non-decreasing left-continuous functions, and it is bounded because
$T_S$ takes values in the finite set $S$. Hence it is the quantile function
of the measure $(T_S\circ Q_\mu)_{\#}\lambda=(T_S)_{\#}(Q_\mu)_{\#}\lambda=(T_S)_{\#}\mu$.
\end{proof}

\begin{theorem}[Best approximation on the domain]\label{thm:projection}
For every $\mu\in\mathcal{P}_2(\R)$ and every $\nu\in\PS$,
\[
  W_2^2\big(\mu,(T_S)_{\#}\mu\big)=\int_\R d(x,S)^2\,d\mu(x)\ \le\ W_2^2(\mu,\nu).
\]
Hence $(T_S)_{\#}\mu$ is a Wasserstein projection of $\mu$ on $\PS$. The
projection is unique if $\mu$ gives no mass to the midpoints
$\{m_1,\dots,m_{K-1}\}$; otherwise the minimizers are exactly the distributions
obtained by splitting the mass of each midpoint $m_k$ in any proportion between
$s_k$ and $s_{k+1}$.
\end{theorem}

\begin{proof}
Let $\nu\in\PS$. Its quantile function $Q_\nu$ takes values in $S$ (because
$\nu(S)=1$), so for every $t\in(0,1)$
\[
  \big(Q_\mu(t)-Q_\nu(t)\big)^2\ \ge\ d\big(Q_\mu(t),S\big)^2
  =\big(Q_\mu(t)-T_S(Q_\mu(t))\big)^2 .
\]
Integrating over $(0,1)$ and using \eqref{eq:W2} and Lemma~\ref{lem:quantile},
\[
  W_2^2(\mu,\nu)\ \ge\ \int_0^1\big(Q_\mu(t)-T_S(Q_\mu(t))\big)^2dt
  = W_2^2\big(\mu,(T_S)_{\#}\mu\big)
  =\int_0^1 d\big(Q_\mu(t),S\big)^2dt=\int_\R d(x,S)^2\,d\mu(x),
\]
the last equality being the change of variables $x=Q_\mu(t)$
($\mu=(Q_\mu)_{\#}\lambda$).

\emph{Uniqueness.} Equality holds iff $Q_\nu(t)$ is a nearest point of
$Q_\mu(t)$ for almost every $t$. If $Q_\mu(t)\notin\{m_k\}$ the nearest point
is unique, so $Q_\nu(t)=T_S(Q_\mu(t))$. The set
$\{t: Q_\mu(t)=m_k\}$ is an interval of length $\mu(\{m_k\})$; on it $Q_\nu$ may
take either value $s_k$ or $s_{k+1}$, and, $Q_\nu$ being non-decreasing, it
equals $s_k$ on an initial sub-interval and $s_{k+1}$ on the rest: the mass
$\mu(\{m_k\})$ is split between $s_k$ and $s_{k+1}$, and every split gives the
same value of the integral.
\end{proof}

\begin{remark}
The proof uses only the pointwise inequality $|x-y|\ge d(x,S)$ for $y\in S$,
so Theorem~\ref{thm:projection} holds for every Wasserstein distance $W_p$,
$p\ge 1$: $(T_S)_{\#}\mu$ minimizes $W_p(\mu,\nu)$ over $\nu\in\PS$, with
$W_p^p(\mu,(T_S)_{\#}\mu)=\int d(x,S)^p\,d\mu$.
\end{remark}

The next result shows that projecting the barycenter is the same as computing
the barycenter directly in the constrained set $\PS$. It relies on the
Hilbert (flat) structure of $(\mathcal{P}_2(\R),W_2)$ and is specific to $p=2$.

\begin{theorem}[Constrained barycenter]\label{thm:constrained}
For every $\nu\in\mathcal{P}_2(\R)$,
\begin{equation}\label{eq:pyth}
  \sum_{i=1}^n w_i W_2^2(\mu_i,\nu)
  =\sum_{i=1}^n w_i W_2^2(\mu_i,\bar\mu)+W_2^2(\bar\mu,\nu).
\end{equation}
Consequently
\[
  \tilde\mu=(T_S)_{\#}\bar\mu\ \in\ \argmin_{\nu\in\PS}\ \sum_{i=1}^n w_i\,W_2^2(\mu_i,\nu),
\]
i.e.\ the simplified barycenter is a barycenter of $\mu_1,\dots,\mu_n$
constrained to the domain $S$, and the set of constrained barycenters coincides
with the set of Wasserstein projections of $\bar\mu$ on $\PS$
(described in Theorem~\ref{thm:projection}).
\end{theorem}

\begin{proof}
Write $Q_i=Q_{\mu_i}$, $\bar Q=Q_{\bar\mu}=\sum_i w_iQ_i$ and $Q=Q_\nu$, and let
$\langle f,g\rangle=\int_0^1 fg\,dt$. By \eqref{eq:W2},
\[
  \sum_i w_i\|Q_i-Q\|^2=\sum_i w_i\|Q_i-\bar Q\|^2+\|\bar Q-Q\|^2
  +2\Big\langle \sum_i w_i(Q_i-\bar Q),\,\bar Q-Q\Big\rangle ,
\]
and the last term vanishes because $\sum_i w_i(Q_i-\bar Q)=\bar Q-\bar Q=0$.
This proves \eqref{eq:pyth}. The first term on the right-hand side of
\eqref{eq:pyth} does not depend on $\nu$, hence minimizing
$\sum_i w_iW_2^2(\mu_i,\nu)$ over $\nu\in\PS$ is equivalent to minimizing
$W_2^2(\bar\mu,\nu)$ over $\nu\in\PS$, whose solutions are given by
Theorem~\ref{thm:projection}.
\end{proof}

\begin{remark}
The set $\PS$ is not convex in the quantile representation (a convex
combination of $S$-valued quantile functions is not $S$-valued), but the proof
does not need convexity: identity \eqref{eq:pyth} is purely algebraic.
Equation \eqref{eq:pyth} also gives the loss due to the constraint:
\[
  \min_{\nu\in\PS}\sum_i w_iW_2^2(\mu_i,\nu)-\min_{\nu\in\mathcal{P}_2(\R)}\sum_i w_iW_2^2(\mu_i,\nu)
  = W_2^2(\bar\mu,\tilde\mu)=\int d(x,S)^2\,d\bar\mu(x).
\]
\end{remark}

\begin{proposition}[Size of the approximation]\label{prop:bounds}
Let $h=\max_k(s_{k+1}-s_k)$ and suppose $\bar\mu([s_1,s_K])=1$ (this holds
when all the $\mu_i$ are supported on $S$, since then $s_1\le Q_i\le s_K$).
Then
\[
  W_2(\bar\mu,\tilde\mu)\le h/2,\qquad
  \big|\mathrm{E}_{\tilde\mu}[X]-\mathrm{E}_{\bar\mu}[X]\big|\le W_1(\bar\mu,\tilde\mu)\le W_2(\bar\mu,\tilde\mu)\le h/2 .
\]
For an integer scale ($h=1$) the simplified barycenter is within $1/2$ from
the exact one, and its mean differs from the exact mean by at most $1/2$.
\end{proposition}

\begin{proof}
For $x\in[s_1,s_K]$, $d(x,S)\le h/2$, hence by Theorem~\ref{thm:projection}
$W_2^2(\bar\mu,\tilde\mu)=\int d(x,S)^2d\bar\mu\le h^2/4$. For the means,
$|\mathrm{E}_{\tilde\mu}X-\mathrm{E}_{\bar\mu}X|=\big|\int_0^1(T_S(\bar Q)-\bar Q)\,dt\big|
\le\int_0^1|T_S(\bar Q)-\bar Q|\,dt=W_1(\bar\mu,\tilde\mu)$, and
$W_1\le W_2$ by Jensen's inequality.
\end{proof}

\begin{remark}[What the simplified barycenter does not preserve]
The projection is optimal with respect to the Wasserstein distance, but it does
not preserve the mean (nor the variance) of the exact barycenter; the error is
bounded by Proposition~\ref{prop:bounds}. If one wants a distribution on $S$
with the same mean as $\bar\mu$, the problem
$\min\{W_2(\bar\mu,\nu):\nu\in\PS,\ \mathrm{E}_\nu X=\mathrm{E}_{\bar\mu}X\}$ is
a different (constrained) problem, whose solution is in general not given by
the nearest-point rule. Other intuitive rules, such as splitting the mass of
each point $x\in(s_k,s_{k+1})$ between $s_k$ and $s_{k+1}$ proportionally to
the distances (which preserves the mean), are not optimal for $W_2$.
\end{remark}

\section{Algorithm}

Given the exact barycenter $\bar\mu=\sum_{l}\pi_l\delta_{x_l}$ (computed from
the averaged quantile functions) and the domain $S$:
\begin{enumerate}
  \item compute the midpoints $m_k=(s_k+s_{k+1})/2$, $k=1,\dots,K-1$;
  \item for $k=1,\dots,K$, set
        $\tilde\pi_k=\sum_{l:\ m_{k-1}<x_l\le m_k}\pi_l$ (with $m_0=-\infty$,
        $m_K=+\infty$);
  \item return $\tilde\mu=\sum_k\tilde\pi_k\delta_{s_k}$.
\end{enumerate}
The cost is $O(L\log K)$, $L$ being the number of support points of
$\bar\mu$. By Theorems~\ref{thm:projection} and~\ref{thm:constrained},
$\tilde\mu$ is both the distribution on $S$ closest to the exact barycenter and
a barycenter of the data constrained to $S$; the convention of assigning the
midpoints to the lower value selects one of the (equivalent) solutions when
$\bar\mu$ has atoms exactly at midpoints.

\end{document}
